% mytest-helpers.tex -- helper macros shared by every regression test.
%
% Input it *after* \documentclass/\usepackage{mypackage}:
%     \input{regression-test}
%     \documentclass{myreport}
%     % mytest-helpers.tex -- helper macros shared by every regression test.
%
% Input it *after* \documentclass/\usepackage{mypackage}:
%     \input{regression-test}
%     \documentclass{myreport}
%     % mytest-helpers.tex -- helper macros shared by every regression test.
%
% Input it *after* \documentclass/\usepackage{mypackage}:
%     \input{regression-test}
%     \documentclass{myreport}
%     % mytest-helpers.tex -- helper macros shared by every regression test.
%
% Input it *after* \documentclass/\usepackage{mypackage}:
%     \input{regression-test}
%     \documentclass{myreport}
%     \input{mytest-helpers}
%
% Every helper writes a single, stable line to the log so that the .tlg
% files stay short and readable.  Helpers use \typeout (not \TYPE) because
% \typeout handles robust/protected commands (babel strings, \textsc, ...).
\makeatletter
% ---------------------------------------------------------------- references
% \SHOWREF{label}: number, page and hyperref anchor stored for a label.
\newcommand\SHOWREF[1]{%
  \typeout{REF #1: number=[\getrefbykeydefault{#1}{}{??}]
    page=[\getrefbykeydefault{#1}{page}{??}]
    anchor=[\getrefbykeydefault{#1}{anchor}{??}]}%
}
% \SHOWTEXT{material}: typeset material in an \hbox and show the box.
\newcommand\SHOWTEXT[1]{\setbox\z@\hbox{#1}\showbox\z@}
% \SHOWSTRING{material}: the *text* of a piece of typeset material,
% obtained by expanding it with \protect = \string-like behaviour.
\newcommand\SHOWSTRING[2][]{%
  \begingroup\let\protect\@unexpandable@protect
  \protected@edef\@tempa{#2}\typeout{STRING#1: \meaning\@tempa}\endgroup}
% ---------------------------------------------------------------- counters
\newcommand\SHOWCOUNTERS{%
  \typeout{COUNTERS: part=\the\c@mypart, chapter=\the\c@mychapter,
    section=\the\c@mysection, subsection=\the\c@mysubsection,
    subsubsection=\the\c@mysubsubsection, figure=\the\c@myfigure,
    subfigure=\the\c@mysubfigure, table=\the\c@mytable,
    program=\the\c@myprogram, equation=\the\c@myequation,
    annex=\the\c@myannex, page=\the\c@page\space(\thepage)}%
}
\newcommand\SHOWTHE{%
  \typeout{THE: part=[\themypart] chapter=[\themychapter]
    section=[\themysection] subsection=[\themysubsection]
    subsubsection=[\themysubsubsection] figure=[\themyfigure]
    subfigure=[\themysubfigure] table=[\themytable]
    program=[\themyprogram] equation=[\themyequation]}%
}
\ExplSyntaxOn
\cs_new:Npn \__mytest_tf:n #1 { \bool_if:nTF {#1} { T } { F } }
\newcommand\SHOWFLAGS{%
  \typeout{FLAGS:~part=\__mytest_tf:n { \mypackage_if_level_used_p:n { part } },~
    chapter=\__mytest_tf:n { \mypackage_if_level_used_p:n { chapter } },~
    section=\__mytest_tf:n { \mypackage_if_level_used_p:n { section } },~
    subsection=\__mytest_tf:n { \mypackage_if_level_used_p:n { subsection } },~
    subsubsection=\__mytest_tf:n { \mypackage_if_level_used_p:n { subsubsection } },~
    annex=\__mytest_tf:n { \l_mypackage_annex_bool },~
    bib=\__mytest_tf:n { \l_mypackage_bibliography_bool },~
    highest=\myhighestsec}%
}
% \SETUSED{<levels>}: marks levels as used, *locally* (inside a \TEST group),
% as if the .aux file had recorded them.
\NewDocumentCommand \SETUSED { m }
  {
    \clist_map_inline:nn {#1}
      {
        \bool_set_true:c { g__mypackage_ ##1 _used_bool }
        \int_set:Nn \g__mypackage_highest_int
          { \int_min:nn { \g__mypackage_highest_int } { \mypackage_level_index:n {##1} } }
      }
  }
% \SETANNEX: annex mode, locally (as inside a myannex environment).
\NewDocumentCommand \SETANNEX { } { \bool_set_true:N \l_mypackage_annex_bool }
% \IFPAGENUMBERING{<style>}{<true>}{<false>}
\NewDocumentCommand \IFPAGENUMBERING { m m m }
  { \mypackage_if_pagenumbering:nTF {#1} {#2} {#3} }
% \SHOWVAR{<expl3 variable name, without backslash>}: its value; integers
% and booleans are evaluated, token lists and sequences shown as strings.
\NewDocumentCommand \SHOWVAR { m }
  {
    \cs_if_exist:cTF {#1}
      {
        \str_case_e:nnF { \str_range:nnn {#1} { -4 } { -1 } }
          {
            { _int } { \typeout { VAR~#1=[ \int_use:c {#1} ] } }
            { bool } { \typeout { VAR~#1=[ \bool_if:cTF {#1} { true } { false } ] } }
            { _seq } { \typeout { VAR~#1=[ \exp_args:Ne \tl_to_str:n { \seq_use:cn {#1} { ~|~ } } ] } }
          }
          { \typeout { VAR~#1=[ \tl_to_str:v {#1} ] } }
      }
      { \typeout { VAR~#1=[UNDEFINED] } }
  }
\ExplSyntaxOff
% ---------------------------------------------------------------- words
\newcommand\SHOWNAMES{%
  \typeout{NAMES-1: contents=[\mycontentsname] lof=[\mylistfigurename]
    lot=[\mylisttablename] bib=[\mybibname] abstract=[\myabstractname]
    chapter=[\mychaptername] part=[\mypartname] figure=[\myfigurename]
    table=[\mytablename]}%
  \typeout{NAMES-2: equation=[\myequationname] acyear=[\myacyearname]
    books=[\mybooksname] articles=[\myarticlesname]
    collections=[\mycollectionsname] manuals=[\mymanualsname]
    theses=[\mythesesname] references=[\myreferencesname]
    reports=[\myreportsname] websites=[\mywebsitesname]
    patents=[\mypatentsname]}%
  \typeout{NAMES-3: program=[\myprogramname] lop=[\mylistprogramname]
    loe=[\mylistequationname] introduction=[\myintroductionname]
    section=[\mysectionname] subsection=[\mysubsectionname]
    subsubsection=[\mysubsubsectionname] annex=[\myannexname]
    see=[\seename]}%
  \typeout{NAMES-4: professor=[\myprofessormname|\myprofessorfname|%
    \myprofessormpname|\myprofessorfpname]
    assistant=[\myassistantmname|\myassistantfname|%
    \myassistantmpname|\myassistantfpname]
    author=[\myauthormname|\myauthorfname|%
    \myauthormpname|\myauthorfpname]}%
}
% ---------------------------------------------------------------- files
% \SHOWLIST{ext}: content of \jobname.ext as left by the *previous* run.
% Must be used before the matching \mytableofcontents/\mylistof... call,
% which truncates the file.
\newcommand\SHOWLIST[1]{\SHOWFILE{\jobname.#1}}
% ---------------------------------------------------------------- option values
% \SHOWOPT\macro: the full expansion of an option macro.
\newcommand\SHOWOPT[1]{%
  \ifdefined#1\protected@edef\@tempa{#1}\else\def\@tempa{UNDEFINED}\fi
  \typeout{OPT \string#1=[\detokenize\expandafter{\@tempa}]}}
\makeatother
\endinput

%
% Every helper writes a single, stable line to the log so that the .tlg
% files stay short and readable.  Helpers use \typeout (not \TYPE) because
% \typeout handles robust/protected commands (babel strings, \textsc, ...).
\makeatletter
% ---------------------------------------------------------------- references
% \SHOWREF{label}: number, page and hyperref anchor stored for a label.
\newcommand\SHOWREF[1]{%
  \typeout{REF #1: number=[\getrefbykeydefault{#1}{}{??}]
    page=[\getrefbykeydefault{#1}{page}{??}]
    anchor=[\getrefbykeydefault{#1}{anchor}{??}]}%
}
% \SHOWTEXT{material}: typeset material in an \hbox and show the box.
\newcommand\SHOWTEXT[1]{\setbox\z@\hbox{#1}\showbox\z@}
% \SHOWSTRING{material}: the *text* of a piece of typeset material,
% obtained by expanding it with \protect = \string-like behaviour.
\newcommand\SHOWSTRING[2][]{%
  \begingroup\let\protect\@unexpandable@protect
  \protected@edef\@tempa{#2}\typeout{STRING#1: \meaning\@tempa}\endgroup}
% ---------------------------------------------------------------- counters
\newcommand\SHOWCOUNTERS{%
  \typeout{COUNTERS: part=\the\c@mypart, chapter=\the\c@mychapter,
    section=\the\c@mysection, subsection=\the\c@mysubsection,
    subsubsection=\the\c@mysubsubsection, figure=\the\c@myfigure,
    subfigure=\the\c@mysubfigure, table=\the\c@mytable,
    program=\the\c@myprogram, equation=\the\c@myequation,
    annex=\the\c@myannex, page=\the\c@page\space(\thepage)}%
}
\newcommand\SHOWTHE{%
  \typeout{THE: part=[\themypart] chapter=[\themychapter]
    section=[\themysection] subsection=[\themysubsection]
    subsubsection=[\themysubsubsection] figure=[\themyfigure]
    subfigure=[\themysubfigure] table=[\themytable]
    program=[\themyprogram] equation=[\themyequation]}%
}
\ExplSyntaxOn
\cs_new:Npn \__mytest_tf:n #1 { \bool_if:nTF {#1} { T } { F } }
\newcommand\SHOWFLAGS{%
  \typeout{FLAGS:~part=\__mytest_tf:n { \mypackage_if_level_used_p:n { part } },~
    chapter=\__mytest_tf:n { \mypackage_if_level_used_p:n { chapter } },~
    section=\__mytest_tf:n { \mypackage_if_level_used_p:n { section } },~
    subsection=\__mytest_tf:n { \mypackage_if_level_used_p:n { subsection } },~
    subsubsection=\__mytest_tf:n { \mypackage_if_level_used_p:n { subsubsection } },~
    annex=\__mytest_tf:n { \l_mypackage_annex_bool },~
    bib=\__mytest_tf:n { \l_mypackage_bibliography_bool },~
    highest=\myhighestsec}%
}
% \SETUSED{<levels>}: marks levels as used, *locally* (inside a \TEST group),
% as if the .aux file had recorded them.
\NewDocumentCommand \SETUSED { m }
  {
    \clist_map_inline:nn {#1}
      {
        \bool_set_true:c { g__mypackage_ ##1 _used_bool }
        \int_set:Nn \g__mypackage_highest_int
          { \int_min:nn { \g__mypackage_highest_int } { \mypackage_level_index:n {##1} } }
      }
  }
% \SETANNEX: annex mode, locally (as inside a myannex environment).
\NewDocumentCommand \SETANNEX { } { \bool_set_true:N \l_mypackage_annex_bool }
% \IFPAGENUMBERING{<style>}{<true>}{<false>}
\NewDocumentCommand \IFPAGENUMBERING { m m m }
  { \mypackage_if_pagenumbering:nTF {#1} {#2} {#3} }
% \SHOWVAR{<expl3 variable name, without backslash>}: its value; integers
% and booleans are evaluated, token lists and sequences shown as strings.
\NewDocumentCommand \SHOWVAR { m }
  {
    \cs_if_exist:cTF {#1}
      {
        \str_case_e:nnF { \str_range:nnn {#1} { -4 } { -1 } }
          {
            { _int } { \typeout { VAR~#1=[ \int_use:c {#1} ] } }
            { bool } { \typeout { VAR~#1=[ \bool_if:cTF {#1} { true } { false } ] } }
            { _seq } { \typeout { VAR~#1=[ \exp_args:Ne \tl_to_str:n { \seq_use:cn {#1} { ~|~ } } ] } }
          }
          { \typeout { VAR~#1=[ \tl_to_str:v {#1} ] } }
      }
      { \typeout { VAR~#1=[UNDEFINED] } }
  }
\ExplSyntaxOff
% ---------------------------------------------------------------- words
\newcommand\SHOWNAMES{%
  \typeout{NAMES-1: contents=[\mycontentsname] lof=[\mylistfigurename]
    lot=[\mylisttablename] bib=[\mybibname] abstract=[\myabstractname]
    chapter=[\mychaptername] part=[\mypartname] figure=[\myfigurename]
    table=[\mytablename]}%
  \typeout{NAMES-2: equation=[\myequationname] acyear=[\myacyearname]
    books=[\mybooksname] articles=[\myarticlesname]
    collections=[\mycollectionsname] manuals=[\mymanualsname]
    theses=[\mythesesname] references=[\myreferencesname]
    reports=[\myreportsname] websites=[\mywebsitesname]
    patents=[\mypatentsname]}%
  \typeout{NAMES-3: program=[\myprogramname] lop=[\mylistprogramname]
    loe=[\mylistequationname] introduction=[\myintroductionname]
    section=[\mysectionname] subsection=[\mysubsectionname]
    subsubsection=[\mysubsubsectionname] annex=[\myannexname]
    see=[\seename]}%
  \typeout{NAMES-4: professor=[\myprofessormname|\myprofessorfname|%
    \myprofessormpname|\myprofessorfpname]
    assistant=[\myassistantmname|\myassistantfname|%
    \myassistantmpname|\myassistantfpname]
    author=[\myauthormname|\myauthorfname|%
    \myauthormpname|\myauthorfpname]}%
}
% ---------------------------------------------------------------- files
% \SHOWLIST{ext}: content of \jobname.ext as left by the *previous* run.
% Must be used before the matching \mytableofcontents/\mylistof... call,
% which truncates the file.
\newcommand\SHOWLIST[1]{\SHOWFILE{\jobname.#1}}
% ---------------------------------------------------------------- option values
% \SHOWOPT\macro: the full expansion of an option macro.
\newcommand\SHOWOPT[1]{%
  \ifdefined#1\protected@edef\@tempa{#1}\else\def\@tempa{UNDEFINED}\fi
  \typeout{OPT \string#1=[\detokenize\expandafter{\@tempa}]}}
\makeatother
\endinput

%
% Every helper writes a single, stable line to the log so that the .tlg
% files stay short and readable.  Helpers use \typeout (not \TYPE) because
% \typeout handles robust/protected commands (babel strings, \textsc, ...).
\makeatletter
% ---------------------------------------------------------------- references
% \SHOWREF{label}: number, page and hyperref anchor stored for a label.
\newcommand\SHOWREF[1]{%
  \typeout{REF #1: number=[\getrefbykeydefault{#1}{}{??}]
    page=[\getrefbykeydefault{#1}{page}{??}]
    anchor=[\getrefbykeydefault{#1}{anchor}{??}]}%
}
% \SHOWTEXT{material}: typeset material in an \hbox and show the box.
\newcommand\SHOWTEXT[1]{\setbox\z@\hbox{#1}\showbox\z@}
% \SHOWSTRING{material}: the *text* of a piece of typeset material,
% obtained by expanding it with \protect = \string-like behaviour.
\newcommand\SHOWSTRING[2][]{%
  \begingroup\let\protect\@unexpandable@protect
  \protected@edef\@tempa{#2}\typeout{STRING#1: \meaning\@tempa}\endgroup}
% ---------------------------------------------------------------- counters
\newcommand\SHOWCOUNTERS{%
  \typeout{COUNTERS: part=\the\c@mypart, chapter=\the\c@mychapter,
    section=\the\c@mysection, subsection=\the\c@mysubsection,
    subsubsection=\the\c@mysubsubsection, figure=\the\c@myfigure,
    subfigure=\the\c@mysubfigure, table=\the\c@mytable,
    program=\the\c@myprogram, equation=\the\c@myequation,
    annex=\the\c@myannex, page=\the\c@page\space(\thepage)}%
}
\newcommand\SHOWTHE{%
  \typeout{THE: part=[\themypart] chapter=[\themychapter]
    section=[\themysection] subsection=[\themysubsection]
    subsubsection=[\themysubsubsection] figure=[\themyfigure]
    subfigure=[\themysubfigure] table=[\themytable]
    program=[\themyprogram] equation=[\themyequation]}%
}
\ExplSyntaxOn
\cs_new:Npn \__mytest_tf:n #1 { \bool_if:nTF {#1} { T } { F } }
\newcommand\SHOWFLAGS{%
  \typeout{FLAGS:~part=\__mytest_tf:n { \mypackage_if_level_used_p:n { part } },~
    chapter=\__mytest_tf:n { \mypackage_if_level_used_p:n { chapter } },~
    section=\__mytest_tf:n { \mypackage_if_level_used_p:n { section } },~
    subsection=\__mytest_tf:n { \mypackage_if_level_used_p:n { subsection } },~
    subsubsection=\__mytest_tf:n { \mypackage_if_level_used_p:n { subsubsection } },~
    annex=\__mytest_tf:n { \l_mypackage_annex_bool },~
    bib=\__mytest_tf:n { \l_mypackage_bibliography_bool },~
    highest=\myhighestsec}%
}
% \SETUSED{<levels>}: marks levels as used, *locally* (inside a \TEST group),
% as if the .aux file had recorded them.
\NewDocumentCommand \SETUSED { m }
  {
    \clist_map_inline:nn {#1}
      {
        \bool_set_true:c { g__mypackage_ ##1 _used_bool }
        \int_set:Nn \g__mypackage_highest_int
          { \int_min:nn { \g__mypackage_highest_int } { \mypackage_level_index:n {##1} } }
      }
  }
% \SETANNEX: annex mode, locally (as inside a myannex environment).
\NewDocumentCommand \SETANNEX { } { \bool_set_true:N \l_mypackage_annex_bool }
% \IFPAGENUMBERING{<style>}{<true>}{<false>}
\NewDocumentCommand \IFPAGENUMBERING { m m m }
  { \mypackage_if_pagenumbering:nTF {#1} {#2} {#3} }
% \SHOWVAR{<expl3 variable name, without backslash>}: its value; integers
% and booleans are evaluated, token lists and sequences shown as strings.
\NewDocumentCommand \SHOWVAR { m }
  {
    \cs_if_exist:cTF {#1}
      {
        \str_case_e:nnF { \str_range:nnn {#1} { -4 } { -1 } }
          {
            { _int } { \typeout { VAR~#1=[ \int_use:c {#1} ] } }
            { bool } { \typeout { VAR~#1=[ \bool_if:cTF {#1} { true } { false } ] } }
            { _seq } { \typeout { VAR~#1=[ \exp_args:Ne \tl_to_str:n { \seq_use:cn {#1} { ~|~ } } ] } }
          }
          { \typeout { VAR~#1=[ \tl_to_str:v {#1} ] } }
      }
      { \typeout { VAR~#1=[UNDEFINED] } }
  }
\ExplSyntaxOff
% ---------------------------------------------------------------- words
\newcommand\SHOWNAMES{%
  \typeout{NAMES-1: contents=[\mycontentsname] lof=[\mylistfigurename]
    lot=[\mylisttablename] bib=[\mybibname] abstract=[\myabstractname]
    chapter=[\mychaptername] part=[\mypartname] figure=[\myfigurename]
    table=[\mytablename]}%
  \typeout{NAMES-2: equation=[\myequationname] acyear=[\myacyearname]
    books=[\mybooksname] articles=[\myarticlesname]
    collections=[\mycollectionsname] manuals=[\mymanualsname]
    theses=[\mythesesname] references=[\myreferencesname]
    reports=[\myreportsname] websites=[\mywebsitesname]
    patents=[\mypatentsname]}%
  \typeout{NAMES-3: program=[\myprogramname] lop=[\mylistprogramname]
    loe=[\mylistequationname] introduction=[\myintroductionname]
    section=[\mysectionname] subsection=[\mysubsectionname]
    subsubsection=[\mysubsubsectionname] annex=[\myannexname]
    see=[\seename]}%
  \typeout{NAMES-4: professor=[\myprofessormname|\myprofessorfname|%
    \myprofessormpname|\myprofessorfpname]
    assistant=[\myassistantmname|\myassistantfname|%
    \myassistantmpname|\myassistantfpname]
    author=[\myauthormname|\myauthorfname|%
    \myauthormpname|\myauthorfpname]}%
}
% ---------------------------------------------------------------- files
% \SHOWLIST{ext}: content of \jobname.ext as left by the *previous* run.
% Must be used before the matching \mytableofcontents/\mylistof... call,
% which truncates the file.
\newcommand\SHOWLIST[1]{\SHOWFILE{\jobname.#1}}
% ---------------------------------------------------------------- option values
% \SHOWOPT\macro: the full expansion of an option macro.
\newcommand\SHOWOPT[1]{%
  \ifdefined#1\protected@edef\@tempa{#1}\else\def\@tempa{UNDEFINED}\fi
  \typeout{OPT \string#1=[\detokenize\expandafter{\@tempa}]}}
\makeatother
\endinput

%
% Every helper writes a single, stable line to the log so that the .tlg
% files stay short and readable.  Helpers use \typeout (not \TYPE) because
% \typeout handles robust/protected commands (babel strings, \textsc, ...).
\makeatletter
% ---------------------------------------------------------------- references
% \SHOWREF{label}: number, page and hyperref anchor stored for a label.
\newcommand\SHOWREF[1]{%
  \typeout{REF #1: number=[\getrefbykeydefault{#1}{}{??}]
    page=[\getrefbykeydefault{#1}{page}{??}]
    anchor=[\getrefbykeydefault{#1}{anchor}{??}]}%
}
% \SHOWTEXT{material}: typeset material in an \hbox and show the box.
\newcommand\SHOWTEXT[1]{\setbox\z@\hbox{#1}\showbox\z@}
% \SHOWSTRING{material}: the *text* of a piece of typeset material,
% obtained by expanding it with \protect = \string-like behaviour.
\newcommand\SHOWSTRING[2][]{%
  \begingroup\let\protect\@unexpandable@protect
  \protected@edef\@tempa{#2}\typeout{STRING#1: \meaning\@tempa}\endgroup}
% ---------------------------------------------------------------- counters
\newcommand\SHOWCOUNTERS{%
  \typeout{COUNTERS: part=\the\c@mypart, chapter=\the\c@mychapter,
    section=\the\c@mysection, subsection=\the\c@mysubsection,
    subsubsection=\the\c@mysubsubsection, figure=\the\c@myfigure,
    subfigure=\the\c@mysubfigure, table=\the\c@mytable,
    program=\the\c@myprogram, equation=\the\c@myequation,
    annex=\the\c@myannex, page=\the\c@page\space(\thepage)}%
}
\newcommand\SHOWTHE{%
  \typeout{THE: part=[\themypart] chapter=[\themychapter]
    section=[\themysection] subsection=[\themysubsection]
    subsubsection=[\themysubsubsection] figure=[\themyfigure]
    subfigure=[\themysubfigure] table=[\themytable]
    program=[\themyprogram] equation=[\themyequation]}%
}
\ExplSyntaxOn
\cs_new:Npn \__mytest_tf:n #1 { \bool_if:nTF {#1} { T } { F } }
\newcommand\SHOWFLAGS{%
  \typeout{FLAGS:~part=\__mytest_tf:n { \mypackage_if_level_used_p:n { part } },~
    chapter=\__mytest_tf:n { \mypackage_if_level_used_p:n { chapter } },~
    section=\__mytest_tf:n { \mypackage_if_level_used_p:n { section } },~
    subsection=\__mytest_tf:n { \mypackage_if_level_used_p:n { subsection } },~
    subsubsection=\__mytest_tf:n { \mypackage_if_level_used_p:n { subsubsection } },~
    annex=\__mytest_tf:n { \l_mypackage_annex_bool },~
    bib=\__mytest_tf:n { \l_mypackage_bibliography_bool },~
    highest=\myhighestsec}%
}
% \SETUSED{<levels>}: marks levels as used, *locally* (inside a \TEST group),
% as if the .aux file had recorded them.
\NewDocumentCommand \SETUSED { m }
  {
    \clist_map_inline:nn {#1}
      {
        \bool_set_true:c { g__mypackage_ ##1 _used_bool }
        \int_set:Nn \g__mypackage_highest_int
          { \int_min:nn { \g__mypackage_highest_int } { \mypackage_level_index:n {##1} } }
      }
  }
% \SETANNEX: annex mode, locally (as inside a myannex environment).
\NewDocumentCommand \SETANNEX { } { \bool_set_true:N \l_mypackage_annex_bool }
% \IFPAGENUMBERING{<style>}{<true>}{<false>}
\NewDocumentCommand \IFPAGENUMBERING { m m m }
  { \mypackage_if_pagenumbering:nTF {#1} {#2} {#3} }
% \SHOWVAR{<expl3 variable name, without backslash>}: its value; integers
% and booleans are evaluated, token lists and sequences shown as strings.
\NewDocumentCommand \SHOWVAR { m }
  {
    \cs_if_exist:cTF {#1}
      {
        \str_case_e:nnF { \str_range:nnn {#1} { -4 } { -1 } }
          {
            { _int } { \typeout { VAR~#1=[ \int_use:c {#1} ] } }
            { bool } { \typeout { VAR~#1=[ \bool_if:cTF {#1} { true } { false } ] } }
            { _seq } { \typeout { VAR~#1=[ \exp_args:Ne \tl_to_str:n { \seq_use:cn {#1} { ~|~ } } ] } }
          }
          { \typeout { VAR~#1=[ \tl_to_str:v {#1} ] } }
      }
      { \typeout { VAR~#1=[UNDEFINED] } }
  }
\ExplSyntaxOff
% ---------------------------------------------------------------- words
\newcommand\SHOWNAMES{%
  \typeout{NAMES-1: contents=[\mycontentsname] lof=[\mylistfigurename]
    lot=[\mylisttablename] bib=[\mybibname] abstract=[\myabstractname]
    chapter=[\mychaptername] part=[\mypartname] figure=[\myfigurename]
    table=[\mytablename]}%
  \typeout{NAMES-2: equation=[\myequationname] acyear=[\myacyearname]
    books=[\mybooksname] articles=[\myarticlesname]
    collections=[\mycollectionsname] manuals=[\mymanualsname]
    theses=[\mythesesname] references=[\myreferencesname]
    reports=[\myreportsname] websites=[\mywebsitesname]
    patents=[\mypatentsname]}%
  \typeout{NAMES-3: program=[\myprogramname] lop=[\mylistprogramname]
    loe=[\mylistequationname] introduction=[\myintroductionname]
    section=[\mysectionname] subsection=[\mysubsectionname]
    subsubsection=[\mysubsubsectionname] annex=[\myannexname]
    see=[\seename]}%
  \typeout{NAMES-4: professor=[\myprofessormname|\myprofessorfname|%
    \myprofessormpname|\myprofessorfpname]
    assistant=[\myassistantmname|\myassistantfname|%
    \myassistantmpname|\myassistantfpname]
    author=[\myauthormname|\myauthorfname|%
    \myauthormpname|\myauthorfpname]}%
}
% ---------------------------------------------------------------- files
% \SHOWLIST{ext}: content of \jobname.ext as left by the *previous* run.
% Must be used before the matching \mytableofcontents/\mylistof... call,
% which truncates the file.
\newcommand\SHOWLIST[1]{\SHOWFILE{\jobname.#1}}
% ---------------------------------------------------------------- option values
% \SHOWOPT\macro: the full expansion of an option macro.
\newcommand\SHOWOPT[1]{%
  \ifdefined#1\protected@edef\@tempa{#1}\else\def\@tempa{UNDEFINED}\fi
  \typeout{OPT \string#1=[\detokenize\expandafter{\@tempa}]}}
\makeatother
\endinput
